%-------------------------------------------------------------------------------
\chapter{Electroweak and top backgrounds}
\label{sec:ewtopbkgs}
%-------------------------------------------------------------------------------

\todo{Describe cross sections and (normalization) uncertainties}

Background contributions from single-boson and diboson production (electroweak background)
and \ttbar and single-top-quark production (top-quark background)
are directly estimated from 
the MC simulated samples. They are dominated by $W \to \tau\nu$
events with $\tau \to \ell\nu\nu$ at low \ut and top-quark pair
production at higher \ut. In the electron channels there is an
additional background contribution from \Wenu events where the reconstructed electron
charge is mismeasured.

%-------------------------------------------------------------------------------
\chapter{Multijet background}
\label{sec:mjbkg}
%-------------------------------------------------------------------------------

Background from QCD multijet (MJ) production cannot be reliably
simulated and has to be derived from data. A short summary of the methodology is described in this paragraph. 
Depending on the lepton
flavour, the MJ background has significant contributions from 
leptons produced in semileptonic decays of heavy quarks, pion and kaon decays,
or photon conversions, and is referred to as fake leptons below. The multijet yields in the electron and muon selections of the
\Wboson-boson analysis are estimated separately for positively and negatively
charged samples. As multijet production is concentrated at lower
values of \ptl, \met\ and \mt\ as compared to the signal, this
background is estimated with a template fit in these kinematic
distributions in a phase-space region with relaxed \met\ and \mt\
selections, called the \textit{fit region}. The \ptl, \met\ and \mt\ templates in the fit region
for the \Wln signal and electroweak and top-quark
backgrounds are obtained using simulation. The multijet templates are
built from events passing the same kinematic selections as the nominal
selection except that the lepton isolation requirement is inverted.
In the final step, the multijet yields in the signal region are
determined by multiplying the fit region yields by a transfer factor
that corrects for the
acceptance of the \met and \mt\ selections and the impact of the inverted isolation
requirements needed to extract the multijet background templates. The multijet distribution is estimated in each channel for the Profile likelihood Fit separately. 

Section~\ref{sec:mj_method} presents the general methodology for the multijet background estimate in $W$ events, and discusses how the various regions and templates are defined and corrected.
Section~\ref{sec:mj_yield} is a summary of the multijet yield derivation, Section~\ref{sec:mj_shape} explains the multijet shape estimation.
Section~\ref{sec:mj_uncertainty} discusses the uncertainties associated with the multijet estimation.
Section~\ref{sec:mj_results} presents the 5 and 13~TeV results in the electron and muon channels. 


%-------------------------------------------------------------------------------
\section{Multijet background methodology}
\label{sec:mj_method}
%------------------------------------------------------------------------------


In this analysis, the $W$‑boson phase space is defined by the following selection criteria: $\ptl > 25~\text{GeV}$; $\met > 25~\text{GeV}$;  $\mt > 50~\text{GeV}$; lepton isolation is imposed using the variable \texttt{ptcone20/Min(pT, 50~GeV)} with the signal region defined by \texttt{ptcone20/Min(pT, 50~GeV)}$ < 0.1$ (the choice of the lepton‑isolation working point is described in Appendix~A.1 of Ref.~\cite{Xu:2657146}); and finally, the hadronic recoil is required to satisfy $u_T < 25~\text{GeV}$.

As shown in Figure~\ref{fig:define4regions}, several event categories are defined for the multijet estimation, corresponding to different regions in phase space and isolation:
%
\begin{itemize}
    \item signal region (SR): isolated leptons (\texttt{ptcone20/Min(pT,~50~GeV)}$<0.1$), signal requirement on \ptl, \met and \mt;
    \item fit region (FR): isolated leptons, relaxed kinematic requirements: $\met>0$~\gev, $m_\text{T}>0$~\gev;
    \item control region 1 (CR1): anti-isolated leptons (\texttt{ptcone20/Min(pT,~50~GeV)}$>0.1$) with FR kinematic requirements;
    \item control region 2 (CR2): anti-isolated leptons with SR kinematic requirements.
\end{itemize}
%
With these definitions, the SR is a subset of the FR, and CR2 is a subset of CR1.
The relaxed kinematic and isolation requirements in FR, CR1 and CR2 are necessary to increase the multijet statistics.

The aim of the procedure is to estimate both the yield and the shape of the multijet background in the signal region (SR).
First, a dedicated recoil correction in CR1 and CR2 needs to be implemented due to the increased activity around the lepton.
The MJ yield is estimated by first determining the MJ shape in the FR through an extrapolation from CR1.
A template fit (fraction fit) is then performed in the FR to obtain the MJ yield within that region.
Finally, the MJ yield in the SR is estimated by using the yield ratio CR2/CR1.
For the MJ shape estimation, the MJ shape in the SR is derived by extrapolating shape information from CR2.
The MJ estimation is done for each fit input observables (\mt, \ptl and \met) in several categories:
%
\begin{itemize}
	\item 8 measurement channels,
	\item 5 $u_T$ categories,
	\item 4 $\eta$ categories.
\end{itemize}
%
\begin{figure}[htbp!]
    \centering
    \subfigure[]{\includegraphics[width=0.45\textwidth]{backgrounds/define4regions1.pdf}} 
    \subfigure[]{\includegraphics[width=0.45\textwidth]{backgrounds/define4regions2.pdf}}
    \caption{Left: The definition of 4 regions used in MJ study according to the division of fiducial space and lepton isolation. Right: Visulatization of the four regions used in MJ study according to the division of fiducial space and lepton isolation.}
    \label{fig:define4regions}
\end{figure}
%

\subsection{Recoil Correction in CR1 and CR2}
\label{sec:mj_recoilCorr}

The recoil reconstruction uses all PFOs in the event, excluding a cone of $\Delta R =0.2$ around the selected leptons. 
To compensate for the underlying event and pile-up energy activity removed by this procedure, another same-size cone of activity in the event is selected, centered at the same $\eta$, but away in azimuth from any lepton and from hard activity; a vector, of magnitude given by the energy observed in this cone and oriented along the removed lepton, is added to the recoil. 
This procedure is successful in the signal region, where leptons are isolated and typically surrounded by soft activity.

As described in Ref~\cite{Xu:2657146} (Section 3.2.1) this cone replacement procedure fails in the anti-isolated regions. 
The compensation can be improved by applying the correction
%
\begin{equation}
  \vec{u}^{\mathrm{corr}} = \vec{u}^{\mathrm{baseline}} + \vec{u}^{\mathrm{iso}} \, , \qquad \mathrm{where}
\end{equation}
%
\begin{equation}
  \vec{u}^{\mathrm{iso}} = k_{\mathrm{iso}} \times \texttt{ptcone20} \vec{n}_l.
\end{equation}
%
where $\vec{u}^{\mathrm{baseline}}$ is the standard recoil reconstruction, $\vec{u}^{\mathrm{iso}}$ represents the present correction, $ \vec{n}_l$ is the unit vector aligned with the lepton direction and $k_{\mathrm{iso}}$ a parameter accounting for a possible difference in scale between the $\texttt{ptcone20}$ energy and the energy removed around the lepton.
Following Ref~\cite{Xu:2657146} the parameters $k_{\mathrm{iso}}$ can be extracted from the data with a linear fit.
This procedure is repeated and really similar values are obtained, the small differences can be explained by small changes in the electron, muon and recoil calibrations.
The estimated values for $k_\mathrm{iso}$ are shown in Table~\ref{tab:kiso}.
%
\begin{table}[h!]
  \centering
  \begin{tabular}{ll}
  \hline
  \textbf{Channel} & \textbf{$k_\mathrm{iso}$} \\
  \hline
  13 TeV  $W^- \rightarrow e^- \nu$  & 1.196 $\pm$ 0.014\\
  13 TeV  $W^+ \rightarrow e^+ \nu$  & 1.178 $\pm$ 0.009\\
  13 TeV  $W^+ \rightarrow mu^+ \nu$  & 1.429 $\pm$ 0.004\\
  13 TeV  $W^- \rightarrow mu^- \nu$  & 1.426 $\pm$ 0.004\\
  \hline
  5 TeV  $W^- \rightarrow e^- \nu$  & 1.141 $\pm$ 0.019\\
  5 TeV  $W^+ \rightarrow e^+ \nu$  & 1.16 $\pm$ 0.02\\
  5 TeV  $W^+ \rightarrow mu^+ \nu$  & 1.422 $\pm$ 0.009\\
  5 TeV  $W^- \rightarrow mu^- \nu$  & 1.383 $\pm$ 0.009\\
  \hline
  \vspace{0.3cm}
  \end{tabular}
  \caption{The fitted values of $k_\mathrm{iso}$ for 5 TeV and 13 TeV, the uncertainties are the statistical uncertainties resulting from the fit.}
  \label{tab:kiso}
\end{table}
%
This recoil correction in the control regions is applied prior to the cut $u_T <25$~GeV.



%-------------------------------------------------------------------------------
\section{Multijet yield estimation}
\label{sec:mj_yield}
%------------------------------------------------------------------------------

The MJ yield estimation is performed in three steps:
%
\begin{enumerate}
  \item Estimate the MJ shape in the FR.
  \item Template Fit to obtain the MJ yield in the FR.
  \item Extrapolate the yield to the SR.
\end{enumerate}
%

\subsection{MJ shape estimation in fit region}
\label{sec:mj_shapeFR}

To determine the number of multijet events in the fit region via a template fit, the shape of the multijet background is determined from CR1.
To remove any dependence of the multijet shape from the isolation criteria itself, control regions CR1 and CR2 are split into successive intervals of the isolation variable as illustrated in Figure~\ref{fig:4regionscut} and are defined as intervals with lepton isolation of \texttt{ptcone20/Min(pT,~50~GeV)} in [0.1, 0.2, 0.3, 0.4, 0.5].
The multijet shape in the FR is then extrapolated from these intervals in the CRs.
%
\begin{figure}[htbp!]
    \centering
    \includegraphics[width=0.45\textwidth]{backgrounds/4regionscut.pdf}
    \caption{Intruducing isolation-dependent regions in CR1 and CR2 to perform an extrapolation to the FR or the SR.}
    \label{fig:4regionscut}
\end{figure}
%

The normalized distribution of multijet events, $H_\text{MJ, CR1}^{i} [x]$ in a given anti-isolated slice (i) is derived for CR1 as 
%
\begin{eqnarray}
	H_\text{MJ, CR1}^{i} [x] &=&  \frac{N_\text{data, CR1}^\text{i} [x] -N_\text{EW, CR1}^\text{i} [x]}{N_\text{total}^{i}} ,
\end{eqnarray}
%
where $N_\text{data, CR1}^\text{i} [x]$ is the number of data events in CR1 for isolation slice $i$ and for the kinematic observable range $x$, $N_\text{EW, CR1}^\text{i} [x]$ is the number of EW background processes and $N_\text{total}^{i}$ is the total number of data events in CR1 for a given isolation slice $i$.
The normalized distribution is a function of $x$, which corresponds to one of the kinematic distributions used in the fit.
% The amount of EW processes in the isolation slices is 3.25\% for the isolation slice [0.1 to 0.2] and 0.37\% for the isolation slice [0.4 to 0.5].
Examples of these distributions for each isolation slice are shown in Figure~\ref{fig:shape_extrapolation}. There is a clear shape dependence of the MJ distribution as a function of the isolation slice. 
%
\begin{figure}[htbp!]
    \centering
    \subfigure[]{\includegraphics[width=0.45\textwidth]{backgrounds/mT_loose_extr_ratio.pdf}} 
    \subfigure[]{\includegraphics[width=0.45\textwidth]{backgrounds/elPt_loose_extr_ratio.pdf}}
    \caption{Visualisation of the isolation dependence of the $m_T$ (left) and $p_T^l$ (right) distributions in the $W^+ \rightarrow e^+\nu$ channel at $\sqrt{s}=13$~TeV. In addition the shape extrapolation to the FR is shown (black line).}
    \label{fig:shape_extrapolation}
\end{figure}
%

To estimate the shape of the multijet background in FR, an averaged two-point extrapolation is used:
%
\begin{equation}
  H_{MJ, FR} [X] = H_{MJ, CR1}^{1}[X] - \Delta^{CR1}[X].
  \label{eqn:FractionFit_Template_Shape_Correction}
\end{equation}
%
where the multijet shape in the first isolation slice ([0.1, 0.2]) is corrected by an extrapolation factor $\Delta^{CR1}[X]$:
%
\begin{eqnarray}
  \Delta^{CR1}[X] = \frac{1}{2} \bigg\{ \frac{H_{MJ, CR1}^{3}[X]-H_{MJ, CR1}^{1}[X]}{2} + \frac{H_{MJ, CR1}^{4}[X]-H_{MJ, CR1}^{2}[X]}{2} \bigg\} .
  \label{eq:deltax}
\end{eqnarray}
%
The results of the extrapolated multijet shape in the FR is also shown in Figure \ref{fig:shape_extrapolation}.
Instead of an averaged two-point extrapolation, a four-point linear extrapolation was also compared.
As seen in Figure \ref{fig:comparison_to_linear}, the difference between the two extrapolation models is negligible. 
%
\begin{figure}[htbp!]
    \centering
    \subfigure[]{\includegraphics[width=0.45\textwidth]{backgrounds/mT_Proper_original_Multiple_bins_loose.pdf}} 
    \subfigure[]{\includegraphics[width=0.45\textwidth]{backgrounds/mT_Proper_original_loose}}
    
    \subfigure[]{\includegraphics[width=0.45\textwidth]{backgrounds/elPt_Proper_original_Multiple_bins_loose.pdf}} 
    \subfigure[]{\includegraphics[width=0.45\textwidth]{backgrounds/elPt_Proper_original_loose}}
    \caption{Comparison of the two point extrapolation with a linear extrapolation to the FR. Left: extrapolation for individual bins of the distributions. The dotted lines are the linear extrapolation, the dashed lines are the two point extrapolation. Right: Control plot showing the similarity of the two extrapolation methods. Results are shown for $m_T$ (top) and $p_T^l$ (bottom).}
    \label{fig:comparison_to_linear}
\end{figure}
%
To estimate the impact of the signal contamination in CR1 for the isolation slice [0.1,0.2] an additional study of the shape extrapolation is performed: the isolation slices are shifted by 0.05, now covering the isolation region [0.15,0.55].
This shift is illustrated in Figure~\ref{fig:comparison_to_shifted} (bottom).
In order to derive the multijet shape in the FR it is necessary to modify Eq.~\ref{eqn:FractionFit_Template_Shape_Correction} by a factor of 1.5 as 
% 
\begin{equation}
  H_{MJ, FR} [X] = H_{MJ, CR1}^{1}[X] - \frac{3}{2} \cdot \Delta^{CR1}[X],
\end{equation}
%
with $H_{MJ, CR1}^{1}[X]$ now being derived in the isolation slice [0.15,0.25].
A comparison between the obtained multijet shapes in the FR is shown in Figure~\ref{fig:comparison_to_shifted} (bottom).
The shape differences are rather small, the effect on the template fit in the fit region has been determined to be negligible compared to the dominating yield uncertainties.
%
\begin{figure}[htbp!]
    \centering
    \subfigure[]{\includegraphics[width=0.45\textwidth]{backgrounds/defineshiftedslices.pdf}}
    
    \subfigure[]{\includegraphics[width=0.45\textwidth]{backgrounds/elPt_Proper_Slice_comp_loose.pdf}} 
    \subfigure[]{\includegraphics[width=0.45\textwidth]{backgrounds/mT_Proper_Slice_comp_loose.pdf}}
    \caption{Top: Visualisation of the shifted isolation regions.
    Bottom: Comparison of the nominal FR extrapolation result to the result obtained with shifted isolation slices. Results are shown for $p_T^l$ (left) and $m_T$ (right).}
    \label{fig:comparison_to_shifted}
\end{figure}
%


\subsection{Template fit in fit region}
\label{sec:mj_fitFR}

The template fit is performed in the FR to obtain the MJ yield in said region.
It is performed as a function of \ptl, \mt or \met, inclusively and in categories of ($\eta$, $u_T$).
Figure~\ref{fig:fit_example} shows the inclusive distributions in FR, after the fit.
The MJ template $H_{MJ}^{FR, [0,0.1]}$ is the one described above, while the EW template, which includes the W and all other EW processes, is derived from the MC.
The fit is done using ROOT's TFractionFitter and returns the fractional change of events for the signal and MJ distributions with respect to the their pre-fit normalizations.
The total number of events in FR, expressed via the fractional changes returned by the fit is
%
\begin{equation}
  H_{Data}^{FR, [0,0.1]} = \alpha \cdot H_{EW}^{FR, [0,0.1]} + (1-\alpha) \cdot H_{MJ}^{FR, [0,0.1]},
  \label{eq:mjfit1}
\end{equation}
%
where the fit parameter $\alpha$ is varied to maximise the agreement.
In this setup $\alpha$ is the electroweak fraction and $(1-\alpha)$ is the multijet fraction.
The MJ yield in the FR is then caluclated as
%
\begin{equation}
  N_{MJ}^{FR} = (1-\alpha) \cdot N_{Data}^{FR}, \qquad \Delta N_{MJ}^{FR} = \Delta\alpha \cdot N_{Data}^{FR}.
\end{equation}
%
During the fit it was ensured that the value of $\alpha$ is always in agreement with one.
As seen in Figure~\ref{fig:fit_example}, the post-fit agreement is generally very good.
For fits performed in categories of ($\eta$, $u_T$) the fits are still in good agreement, due to the limited statistics the uncertainty on the fit parameter $\alpha$ increases.
As an example, Figure \ref{fig:fit_example2} shows the result for ($\eta_2$, $u_{T,2}$).

As an additional cross check the fit is also implemented with two free parameters:
%
\begin{equation}
  H_{Data}^{FR, [0,0.1]} =  \alpha_{EW} \cdot H_{EW}^{FR, [0,0.1]} + \alpha_{MJ} \cdot H_{MJ}^{FR, [0,0.1]}.
\end{equation}
%
This changes the result for the MJ yield in the FR by less than one event in each fit.
Hence, this source of uncertainty is also considered negligible.
%
\begin{figure}[htbp!]
    \centering
    \subfigure[]{\includegraphics[width=0.45\textwidth]{backgrounds/FR_ptl.pdf}}
    
    \subfigure[]{\includegraphics[width=0.45\textwidth]{backgrounds/FR_mT.pdf}} 
    \subfigure[]{\includegraphics[width=0.45\textwidth]{backgrounds/FR_met.pdf}}
    \caption{Results for the template fit in the FR.}
    \label{fig:fit_example}
\end{figure}
%
\begin{figure}[htbp!]
    \centering
    \subfigure[]{\includegraphics[width=0.45\textwidth]{backgrounds/FR_ptl_ut2_eta2.pdf}}
    
    \subfigure[]{\includegraphics[width=0.45\textwidth]{backgrounds/FR_mT_ut2_eta2.pdf}} 
    \subfigure[]{\includegraphics[width=0.45\textwidth]{backgrounds/FR_met_ut2_eta2.pdf}}
    \caption{Results for the template fit in the FR in the category ($\eta_2$, $u_{T,2}$).}
    \label{fig:fit_example2}
\end{figure}
%

\subsection{Estimate of MJ yield in the signal region}
\label{sec:mj_resultSR}

The fitted multijet yield in the FR is extrapolated to the SR. 
For this the ratio of MJ events between CR1 and CR2 is used, however, this ratio is also dependent of the isolation slice.
For each isolation slice the corrected MJ yield is calculated as 
%
\begin{equation}
  N_{MJ}([a,b]) = N_{MJ}^{FR} \cdot \frac{N_{MJ}^{CR2,[a,b]}}{N_{MJ}^{CR1[a,b]}} = N_{MJ}^{FR} \cdot \frac{N_\mathrm{MJ-CR2}^i}{N_\mathrm{MJ-CR1}^i}.
  \label{eq:mjfit2}
\end{equation}
%
Here $N_{MJ}([a,b])$ is defined as the number of MJ events in the SR, estimated in the anti-isolated slice ($i$).
The uncertainty on this estimation is calucalted as
%
\begin{equation}
  \Delta N_{MJ}([a,b] = \sqrt{\left(\Delta N_{MJ}^{FR} \cdot \frac{N_{MJ}^{CR2,[a,b]}}{N_{MJ}^{CR1,[a,b]}} \right)^2+ \left( N_{MJ}^{FR} \cdot \Delta\left( \frac{N_{MJ}^{CR2,[a,b]}}{N_{MJ}^{CR1,[a,b]}}\right)\right)^2}.
\end{equation} 
%
The uncertainty $\Delta N_{MJ}^{FR}$ is propagated from the template fit in the last step.
The second part of the equation accounts for statistical uncertainties of the number of events in CR1 and CR2.

The values of $N_{MJ}([a,b])$ as a function of the isolation slice is shown in Figure~\ref{fig:inclusive_extrapolation} for \ptl, \met\ and \mt in the fully inclusive case.
Figure~\ref{fig:categorised_extrapolation} shows the same for \ptl, \met\ and \mt in the category ($\eta_2$, $u_{T,2}$).
%
\begin{figure}[htbp!]
    \centering
    \subfigure[]{\includegraphics[width=0.45\textwidth]{backgrounds/Extrap_met.pdf}}
    
    \subfigure[]{\includegraphics[width=0.45\textwidth]{backgrounds/Extrap_mT.pdf}} 
    \subfigure[]{\includegraphics[width=0.45\textwidth]{backgrounds/Extrap_ptl.pdf}}
    \caption{Extrapolation of the multijet yield to the SR for observables \met (top) \mt (left) and \ptl (right). The extrapolation target is 0.025 in both cases.
    Several possible extrapolations are shown: a linear extrapolation based on all isolation slices (blue), a quadratic extrapolation (red) and a linear extrapolation only including the first three isolation slices (green).
    The uncertainties on the extrapolation are visulaised for the four-point-linear case (blue suqare) as well as the quadratic extrapolation (red square).}
    \label{fig:inclusive_extrapolation}
\end{figure}
%
\begin{figure}[htbp!]
    \centering
    \subfigure[]{\includegraphics[width=0.45\textwidth]{backgrounds/Extrap_ut2eta2_met.pdf}}
    
    \subfigure[]{\includegraphics[width=0.45\textwidth]{backgrounds/Extrap_ut2eta2_mT.pdf}} 
    \subfigure[]{\includegraphics[width=0.45\textwidth]{backgrounds/Extrap_ut2eta2_ptl.pdf}}
    \caption{Extrapolation of the multijet yield to the SR for observables \met (top) \mt (left) and \ptl (right) in the ($\eta_2$, $u_{T,2}$) category. The extrapolation target is 0.025 in both cases.
    The extrapolation target is 0.025 in both cases.
    Several possible extrapolations are shown: a linear extrapolation based on all isolation slices (blue), a quadratic extrapolation (red) and a linear extrapolation only including the first three isolation slices (green).
    The uncertainties on the extrapolation are visulaised for the four-point-linear case (blue suqare) as well as the quadratic extrapolation (red square).}
    \label{fig:categorised_extrapolation}
\end{figure}
%

The four isolation slices are fit to a linear function and the number of MJ events in the SR is evaluated at an isolation value of $0.025$ \,.
This average isolation of the multijet background has been studied in Ref~\cite{Xu:2657146} (Section 3.2.3) with $b\overline{b}$ and $c\overline{c}$ samples in the muon and antimuon channel.
The fassociated yield uncertainty in the SR is given by the linear fit uncertainty evaluated at the isolation value of $0.025$ \,.
The difference to an extrapolation to zero is always covered with the yield uncertainty.


\subsection{Cross Checks and Closure}

The default method for the MJ yield extrapolation is a linear fit that uses all four isolation slices.
The associated uncertainty on the MJ yield is taken from the uncertainty of this linear fit.
In addition to the default method, two alternative approaches were investigated.
The first is a three-point linear fit, shown in green, in which the fourth isolation point is excluded.
The second is a quadratic fit, shown in red.

For each observable, $160$ extrapolations are performed across the ($\eta, u_T$) fit categories: 8 (fit channels) $\cdot$ 4($\eta$) $\cdot$ 5($u_T$)
A comparison between the four-point linear fit and the three-point linear fit shows disagreement in only 2 out of 160 extrapolations, i.e.\ the difference between the approaches is not covered by the four-point linear fit uncertainty.
In these cases the associated yield uncertainty was inflated by using the difference between the approaches as the new uncertainty.

A comparison between the four-point linear fit and the quadratic fit shows disagreement in 19 out of 160 extrapolations, some examples are shown in Figure~\ref{fig:mj_quad}.
These cases show that the quadratic fit is either strongly driven by the fourth isolation point or is overfitting the data.
Based on these studies, the default approach for the ($\eta, u_T$) fit categories is to use the linear fit with all four isolation slices and to assign the MJ yield uncertainty from the uncertainty of this fit.
The quadratic fit and the three-point linear fit are therefore not retained as part of the default procedure.
%
\begin{figure}[htbp!]
    \centering
    \includegraphics[width=\textwidth]{backgrounds/mj_quadExp.pdf}
    \caption{Examples for extrapolations where a quadratic assumption does not hold. The quadratic fit (red line) is either driven by the fourth point or overfitting, leading to unphysical results for the multijet yield.}
    \label{fig:mj_quad}
\end{figure}
%


The slopes of the extrapolation are visualised in Figure~\ref{fig:mj_slopes}.
A clear dependence on both $\eta$ and $u_T$ is visible.
An extrapolation in the ($\eta, u_T$)-inclusive case does not capture this dependence.
%
\begin{figure}[htbp!]
    \centering
    \includegraphics[width=\textwidth]{backgrounds/mj_slopes.pdf}
    \caption{Dependence of the multijet slope with respect to the ($\eta, u_T$) category that it was obtained in. It is clearly visible that the slope depends strongly on $u_T$. A $u_T$-inclusive model does not account for this dependency.}
    \label{fig:mj_slopes}
\end{figure}
%

The inclusive case (integrated over $\eta$ and $u_T$) is shown in Figure~\ref{fig:inclusive_extrapolation}.
In each observable there is a significant difference between the quadratic fit extrapolation and the linear fit extrapolation.
In order to address this a dedicated closure check is performed for each observable (\ptl, \mt, \met): the sum of the yield results in ($\eta, u_T$) categories is compared to the inclusive yield, estimated by a linear or a quadratic fit.
As an example, for \ptl the comparison can be written as:
%
\begin{equation}
  MJ(p_T^l) = \sum_\eta MJ(p_T^l(\eta)) = \sum_\eta \sum_{u_T} MJ(p_T^l(\eta,u_T)).
  \label{eq:mj_closure}
\end{equation}
%
Figure~\ref{fig:mj_closure} shows the difference between the inclusive yield and the sum of the $u_T$ categorised yields, both fully inclusive and as a function of $\eta$.
The blue dots visualise the results for a full linear extrapolation, closure with 0 is always obtained.
The red dots result from choosing a quadratic extrapolation for the inclusive model, which leads to significant differences to the sum of the $u_T$ categories.
These studies again show that the linear extrapolation leads to consistent results, while the quadratic extrapolation suffers from systematic biases.
Therefore, no improved model for the multijet yield estimation in the inclusive case is required.
%
\begin{figure}[htbp!]
    \centering
    \includegraphics[width=\textwidth]{backgrounds/mj_closure.pdf}
    \caption{Closure cross-check for the MJ yield estimation. The sum of the $u_T$-differential yields is compared with the yield estimated with the inclusive distribution. The blue points visualise the closure assuming a linear extrapolation in the inclusive case, the red points are obtained with a quadratic extrapolation.}
    \label{fig:mj_closure}
\end{figure}
%

Finally, the multijet yields are compared between positive and negative charges, since the multijet yield is charge-independent.
The difference for each channel is shown in Figure~\ref{fig:mj_charge_dependence}.
Positive- and negative-charge distributions agree well, especially in the ($\eta, u_T$) categories.
In the inclusive case, the positive-charge multijet yields are consistently slightly higher, which again indicates the incompleteness in the model.
%
\begin{figure}[htbp!]
    \centering
    \includegraphics[width=\textwidth]{backgrounds/mj_charge_dependence.pdf}
    \caption{Difference between multijet yields in the positive and negative charged electron channels. In the ($\eta, u_T$) categories the results are in agreement with zero, in the inclusive cases there is a slight bias towards the positive charged channel.}
    \label{fig:mj_charge_dependence}
\end{figure}
%

%-------------------------------------------------------------------------------
\section{Multijet shape estimation}
\label{sec:mj_shape}
%------------------------------------------------------------------------------


In the fully inclusive and $u_T$-inclusive categories the shape of the multijet background in the signal region is determined by extrapolating the distributions observed in different anti-isolated slices in CR2.
Four isolation ranges (\texttt{ptcone20/Min(pT, 50 GeV)} within [0.1, 0.2, 0.3, 0.4, 0.5] as above) are used to determine the extrapolation.
For each observable $X$ ($X = p_\text{T}^\ell$, \met, $m_\text{T}$ and so on), the shape extrapolation uses a formula similar to Eq.~\ref{eq:deltax}, except the MJ shapes in CR1 are replaced by those in CR2:
%
\begin{equation}
  H_{MJ}^{SR}[X] = H_{MJ}^{[0.1,0.2]}[X] + \Delta^{CR2}[X] .
\label{eq:CR2_shape}
\end{equation}
%
In the ($\eta, u_T$) categories such an extrapolation is not applied, because the isolation slices are heavily affected by statistical fluctuations.
Instead, the isolation slice in CR2 which is closest to the SR ([0.1,0.2]) is used as a shape in the SR.

The resulting SR shape is affected by statistical fluctuations due to the limited number of events in the underlying CR2 slices.
These fluctuations may lead to unphysical fluctuations in the multijet template.
To suppress this effect, a KDE (Kernel density estimation) smoothing algorithm is applied to the extracted multijet distribution.
In this algorithm, the bandwidth parameter $\rho$ controls the degree of smoothing. 
While larger values of $\rho$ reduce bin-to-bin fluctuations more effectively, they can wash out features of the underlying distribution.
Conversely, smaller values of $\rho$ preserve more of the intrinsic structure at the expense of retaining additional statistical noise.
Since there is no single correct choice for the optimal value of $\rho$, the smoothing is repeated for several values of $\rho$.
The resulting variation in the smoothed multijet template, relative to a nominal smoothing choice, is treated as the systematic shape uncertainty.
Examples for the estimated multijet shape are shown in Figure~\ref{fig:mj_shape}.
%
\begin{figure}[htbp!]
    \centering
    \includegraphics[width=\textwidth]{backgrounds/mj_shape_variation.pdf}
    \caption{Examples for multijet shape variations for different choices of the parameter $\rho$. 
    The red line shows the result for $\rho=0.75$, the default value for the smoothing parameter. 
    The systematic uncertainty is determined by using $\rho=0.5$ and $\rho=1$.}
    \label{fig:mj_shape}
\end{figure}
%


%-------------------------------------------------------------------------------
\section{Systematic Uncertainties}
\label{sec:mj_uncertainty}
%-------------------------------------------------------------------------------

The uncertainty in the MJ background estimation dominantly arises from 3 aspects: systematic uncertainties in the estimation of the yield, systematics in the determination of the shape, and statistics.
These three components have different correlation patterns: systematic uncertainties in yield and shape tend to be correlated across the channels, while statistical uncertainty is uncorrelated due to its statistical nature.
All sources of uncertainty investigated in this analysis are discussed below.

The following sources of uncertainty are assumed uncorrelated among channels:

%
\begin{itemize}
  \item Bin-by-bin uncorrelated data and MC statistics: the statistical uncertainty from the limited data and MC samples is taken from the results of shape extrapolation according to Eq.~\ref{eq:CR2_shape}.
  \item Uncertainty in the yield caused by the linear extrapolation as a function of isolation.
  The uncertainty of the linear extrapolation comes from the error bars in Figure~\ref{fig:categorised_extrapolation} that account for the statistical uncertainty in the control regions and the fraction uncertainty of MJ in the FR.
  This source of uncertainty is the dominating yield uncertainty and propagated to the final output of MJ estimation.
  \item Uncertainty of the shape due to the choice of the smoothing parameter. 
  The result of the smoothing is impacted by the amount of statistical fluctuations in the unsmoothed distribution.
  This source of uncertainty is propagated to the final output of MJ estimation.
\end{itemize}
%
The following potential sources of uncertainty are correlated across channels, and across bins in a given kinematic distribution:
%
\begin{itemize}
  \item Uncertainty due to the shape correction which is used to estimate the shape in the SR or FR.
  A comparison to the shape estimated with moved slices ([0.15, 0.55]) and with a full linear extrapolation has been implemented.
  The differences were found to be negligible.
  \item Uncertainty due to the chosen number of parameters in the template fit.
  A two parameter fit was implemented as a comparison, the difference is less than one event and hence negligible.
  \item Uncertainty in the luminosity and in the cross sections of the simulated samples, assumed in the subtraction of the EW+top processes from CR1 and CR2.
  Negligible with respect to the other uncertainties.
  \item Uncertainty in the choice of the extrapolation target for the isolation scan, estimated by changing the isolation target from 0.025 to 0.
  This difference is always covered by the uncertainty of the linear extrapolation.
  \item The uncertainty in the isolation correction ($k_{\text{iso}}$ in Table~\ref{tab:kiso}) is estimated by changing the isolation correction by 10\%.
  The effect on the multijet yield and shape is negligible.
  \item The uncertainty of the yield and shape due to the dependency of the determination on the assumed value for $\Delta m_W$.
  This impact is further discussed in appendix~\ref{a:mj_mW_dependency}.
\end{itemize}
%
Among these sources the uncertainty due to the linear extrapolation is dominant for the multijet yield, while the multijet shape uncertainty is dominated by the choice of the smoothing parameter $\rho$.
Other sources are found to be sub-dominant in comparison.


%-------------------------------------------------------------------------------
\section{Yield Results}
\label{sec:mj_results}
%-------------------------------------------------------------------------------

The results of for the multijet yields can be found in Tables~\ref{tab:wplusmunu5TeV}, \ref{tab:wminusmunu5TeV}, \ref{tab:wplusenu5TeV}, \ref{tab:wminusenu5TeV}, \ref{tab:wplusmunu13TeV}, \ref{tab:wminusmunu13TeV}, \ref{tab:wplusenu13TeV}, \ref{tab:wminusenu13TeV}.
The results in each individual category agrees within uncertainties for the three observables (\ptl, \mt, \met), with only few exceptions.
The average multijet yield and the largest of the three uncertainties is propagated to the profile likelihood fit of $m_W$.

\begin{table}[htbp]
\centering
\scriptsize
\caption{SR MJ yield for wplusmunu5TeV}
\label{tab:wplusmunu5TeV}
\setlength{\tabcolsep}{6pt}
\begin{tabular}{lrrrrr}
\hline
bin & ptl & mT & ptn & Inclusive & Consistency \\
\hline
uT1\_lEta1 & 23.7247 $\pm$ 91.2 & 5.59834 $\pm$ 36.8 & 40.4801 $\pm$ 33.1 &  & $\checkmark$ \\
uT1\_lEta2 & 7.51326 $\pm$ 57.3 & 9.29088 $\pm$ 31.8 & 25.2978 $\pm$ 32.5 &  & $\checkmark$ \\
uT1\_lEta3 & -0.00383026 $\pm$ 23.3 & -6.65903e-05 $\pm$ 42.6 & -1.20331e-07 $\pm$ 18.7 &  & $\checkmark$ \\
uT1\_lEta4 & 1.52105e-06 $\pm$ 23.4 & 2.55343 $\pm$ 22.7 & 3.63404 $\pm$ 17.5 &  & $\checkmark$ \\
uT2\_lEta1 & 93.2538 $\pm$ 60.2 & 113.29 $\pm$ 38.3 & 109.177 $\pm$ 37.3 &  & $\checkmark$ \\
uT2\_lEta2 & 0.378053 $\pm$ 451 & 24.6706 $\pm$ 25 & 22.1125 $\pm$ 24.9 &  & $\checkmark$ \\
uT2\_lEta3 & 0.000470396 $\pm$ 17.2 & 43.2246 $\pm$ 26.8 & 48.0363 $\pm$ 27.8 &  & x \\
uT2\_lEta4 & 4.39662e-07 $\pm$ 8 & 1.07632 $\pm$ 13 & 0.201619 $\pm$ 2.67 &  & $\checkmark$ \\
uT3\_lEta1 & 42.3747 $\pm$ 27 & 49.3388 $\pm$ 28.1 & 49.4349 $\pm$ 28.2 &  & $\checkmark$ \\
uT3\_lEta2 & 1.06314 $\pm$ 27.5 & -1.51688 $\pm$ 23 & -1.44167 $\pm$ 22.8 &  & $\checkmark$ \\
uT3\_lEta3 & -0.635884 $\pm$ 25.7 & -1.7703 $\pm$ 25.9 & -1.79297 $\pm$ 26.3 &  & $\checkmark$ \\
uT3\_lEta4 & 2.69669 $\pm$ 13 & 1.81545 $\pm$ 16.4 & 1.9409 $\pm$ 16.2 &  & $\checkmark$ \\
uT4\_lEta1 & 20.1046 $\pm$ 28.3 & 19.6703 $\pm$ 28.2 & 19.6892 $\pm$ 28.2 &  & $\checkmark$ \\
uT4\_lEta2 & 27.731 $\pm$ 23.8 & 27.6273 $\pm$ 23.6 & 27.7418 $\pm$ 23.6 &  & $\checkmark$ \\
uT4\_lEta3 & 35.4439 $\pm$ 23.3 & 39.438 $\pm$ 25.6 & 38.9707 $\pm$ 25.3 &  & $\checkmark$ \\
uT4\_lEta4 & -10.7584 $\pm$ 16.8 & -10.7603 $\pm$ 16.9 & -10.6604 $\pm$ 16.7 &  & $\checkmark$ \\
uT5\_lEta1 & 48.6283 $\pm$ 24.8 & 47.5385 $\pm$ 24.1 & 47.2607 $\pm$ 24 &  & $\checkmark$ \\
uT5\_lEta2 & 12.8897 $\pm$ 19.1 & 13.0165 $\pm$ 19.3 & 12.9764 $\pm$ 19.2 &  & $\checkmark$ \\
uT5\_lEta3 & 14.0618 $\pm$ 23.4 & 13.8743 $\pm$ 23 & 13.9082 $\pm$ 23.1 &  & $\checkmark$ \\
uT5\_lEta4 & -5.55463 $\pm$ 17.5 & -5.23126 $\pm$ 16.3 & -5.27001 $\pm$ 16.4 &  & $\checkmark$ \\
\hline
lEta1 & 262.999 $\pm$ 69.4 & 262.081 $\pm$ 65.3 & 259.397 $\pm$ 64.6 &  & $\checkmark$ \\
lEta2 & 60.7296 $\pm$ 55 & 53.5828 $\pm$ 52.7 & 53.3966 $\pm$ 52.4 &  & $\checkmark$ \\
lEta3 & 111.967 $\pm$ 56.4 & 133.96 $\pm$ 61.5 & 132.692 $\pm$ 60.9 &  & $\checkmark$ \\
lEta4 & -8.64699 $\pm$ 34.7 & -19.4246 $\pm$ 42 & -19.4332 $\pm$ 42.2 &  & $\checkmark$ \\
\hline
 &  &  &  & mT=427.573 $\pm$ 111 &  \\
 &  &  &  & muPt=415.242 $\pm$ 112 &  \\
 &  &  &  & ptn=423.889 $\pm$ 110 &  \\
\hline
\end{tabular}
\end{table}
\begin{table}[htbp]
\centering
\scriptsize
\caption{SR MJ yield for wminusmunu5TeV}
\label{tab:wminusmunu5TeV}
\setlength{\tabcolsep}{6pt}
\begin{tabular}{lrrrrr}
\hline
bin & ptl & mT & ptn & Inclusive & Consistency \\
\hline
uT1\_lEta1 & 1.4033 $\pm$ 14 & -0.000126028 $\pm$ 12.6 & 0.0356501 $\pm$ 19.2 &  & $\checkmark$ \\
uT1\_lEta2 & 35.381 $\pm$ 51.3 & 5.71757 $\pm$ 23.3 & 17.7264 $\pm$ 23.3 &  & $\checkmark$ \\
uT1\_lEta3 & -0.00255146 $\pm$ 127 & 35.4835 $\pm$ 37.8 & 32.7865 $\pm$ 29.8 &  & $\checkmark$ \\
uT1\_lEta4 & 13.6104 $\pm$ 37.2 & 5.39784e-05 $\pm$ 6.91 & 0.982089 $\pm$ 17.9 &  & $\checkmark$ \\
uT2\_lEta1 & 65.4032 $\pm$ 35.6 & 61.1853 $\pm$ 24.3 & 64.4133 $\pm$ 25.3 &  & $\checkmark$ \\
uT2\_lEta2 & 69.0405 $\pm$ 37.3 & 49.1473 $\pm$ 23.5 & 50.6647 $\pm$ 23.6 &  & $\checkmark$ \\
uT2\_lEta3 & 95.6095 $\pm$ 46.4 & 63.2507 $\pm$ 28.8 & 70.7869 $\pm$ 30 &  & $\checkmark$ \\
uT2\_lEta4 & 6.09635 $\pm$ 18.5 & 5.83996 $\pm$ 14.2 & 3.45946 $\pm$ 10.9 &  & $\checkmark$ \\
uT3\_lEta1 & 53.6295 $\pm$ 26.2 & 51.353 $\pm$ 24 & 52.4434 $\pm$ 24.5 &  & $\checkmark$ \\
uT3\_lEta2 & 28.0297 $\pm$ 20.3 & 31.6621 $\pm$ 21.4 & 31.5589 $\pm$ 21.4 &  & $\checkmark$ \\
uT3\_lEta3 & 31.5327 $\pm$ 23 & 30.6101 $\pm$ 21.5 & 28.788 $\pm$ 20.2 &  & $\checkmark$ \\
uT3\_lEta4 & -5.25606 $\pm$ 13.4 & -9.3438 $\pm$ 16.4 & -9.30839 $\pm$ 16.5 &  & $\checkmark$ \\
uT4\_lEta1 & 4.31223 $\pm$ 24.2 & 4.13702 $\pm$ 24.1 & 4.10295 $\pm$ 23.9 &  & $\checkmark$ \\
uT4\_lEta2 & -3.33207 $\pm$ 18.7 & -3.48685 $\pm$ 19.1 & -3.53792 $\pm$ 19.3 &  & $\checkmark$ \\
uT4\_lEta3 & 23.4356 $\pm$ 22 & 23.4116 $\pm$ 21.8 & 22.9954 $\pm$ 21.4 &  & $\checkmark$ \\
uT4\_lEta4 & -20.2232 $\pm$ 14.8 & -20.0797 $\pm$ 14.4 & -20.0413 $\pm$ 14.4 &  & $\checkmark$ \\
uT5\_lEta1 & 30.5411 $\pm$ 20.5 & 30.6344 $\pm$ 20.5 & 30.4688 $\pm$ 20.4 &  & $\checkmark$ \\
uT5\_lEta2 & 7.17811 $\pm$ 17.4 & 6.92402 $\pm$ 16.8 & 6.95884 $\pm$ 16.9 &  & $\checkmark$ \\
uT5\_lEta3 & -19.2188 $\pm$ 15.5 & -19.2139 $\pm$ 15.5 & -19.3369 $\pm$ 15.6 &  & $\checkmark$ \\
uT5\_lEta4 & -5.13067 $\pm$ 11.5 & -5.37682 $\pm$ 12.1 & -5.42417 $\pm$ 12.2 &  & $\checkmark$ \\
\hline
lEta1 & 144.023 $\pm$ 57.7 & 140.595 $\pm$ 55.3 & 139.879 $\pm$ 55.1 &  & $\checkmark$ \\
lEta2 & 71.9862 $\pm$ 48.6 & 68.9123 $\pm$ 45.8 & 68.572 $\pm$ 45.6 &  & $\checkmark$ \\
lEta3 & 134.528 $\pm$ 53.7 & 130.506 $\pm$ 50.2 & 128.687 $\pm$ 49.6 &  & $\checkmark$ \\
lEta4 & -25.2693 $\pm$ 31.3 & -34.1857 $\pm$ 36 & -33.4085 $\pm$ 35.3 &  & $\checkmark$ \\
\hline
 &  &  &  & mT=309.652 $\pm$ 93.7 &  \\
 &  &  &  & muPt=317.142 $\pm$ 97.7 &  \\
 &  &  &  & ptn=306.68 $\pm$ 92.8 &  \\
\hline
\end{tabular}
\end{table}
\begin{table}[htbp]
\centering
\scriptsize
\caption{SR MJ yield for wplusenu5TeV}
\label{tab:wplusenu5TeV}
\setlength{\tabcolsep}{6pt}
\begin{tabular}{lrrrrr}
\hline
bin & ptl & mT & ptn & Inclusive & Consistency \\
\hline
uT1\_lEta1 & 24.2523 $\pm$ 98.7 & 49.0949 $\pm$ 50.3 & 11.5766 $\pm$ 36.3 &  & $\checkmark$ \\
uT1\_lEta2 & 0.00086899 $\pm$ 110 & 41.7575 $\pm$ 43.4 & 31.1482 $\pm$ 34 &  & $\checkmark$ \\
uT1\_lEta3 & 40.1963 $\pm$ 75.9 & 47.7975 $\pm$ 38.1 & 46.2107 $\pm$ 33.9 &  & $\checkmark$ \\
uT1\_lEta4 & 5.4981e-05 $\pm$ 46.2 & 7.71528 $\pm$ 26 & 26.5713 $\pm$ 32.7 &  & $\checkmark$ \\
uT2\_lEta1 & 128.592 $\pm$ 49.4 & 159.372 $\pm$ 45.1 & 173.914 $\pm$ 48.6 &  & $\checkmark$ \\
uT2\_lEta2 & 62.7056 $\pm$ 43.2 & 82.9394 $\pm$ 43.6 & 79.1855 $\pm$ 41.9 &  & $\checkmark$ \\
uT2\_lEta3 & 113.635 $\pm$ 48.4 & 148.946 $\pm$ 54.5 & 149.601 $\pm$ 54.7 &  & $\checkmark$ \\
uT2\_lEta4 & 235.205 $\pm$ 94.9 & 253.418 $\pm$ 93.6 & 243.398 $\pm$ 89.9 &  & $\checkmark$ \\
uT3\_lEta1 & 140.265 $\pm$ 53.2 & 135.655 $\pm$ 50.5 & 134.691 $\pm$ 50.2 &  & $\checkmark$ \\
uT3\_lEta2 & 53.645 $\pm$ 50.5 & 52.8788 $\pm$ 49.3 & 52.3818 $\pm$ 48.8 &  & $\checkmark$ \\
uT3\_lEta3 & 130.636 $\pm$ 63.3 & 131.902 $\pm$ 63.5 & 134.082 $\pm$ 64.5 &  & $\checkmark$ \\
uT3\_lEta4 & 151.363 $\pm$ 84.1 & 155.209 $\pm$ 86.1 & 155.513 $\pm$ 86.3 &  & $\checkmark$ \\
uT4\_lEta1 & 125.826 $\pm$ 48.9 & 129.28 $\pm$ 49.9 & 130.974 $\pm$ 50.6 &  & $\checkmark$ \\
uT4\_lEta2 & 164.513 $\pm$ 48.8 & 163.448 $\pm$ 48.4 & 163.127 $\pm$ 48.3 &  & $\checkmark$ \\
uT4\_lEta3 & 145.334 $\pm$ 60.9 & 148.104 $\pm$ 62 & 149.862 $\pm$ 62.7 &  & $\checkmark$ \\
uT4\_lEta4 & 55.114 $\pm$ 97.6 & 53.8644 $\pm$ 95.4 & 54.0997 $\pm$ 95.8 &  & $\checkmark$ \\
uT5\_lEta1 & 72.9901 $\pm$ 44.8 & 71.8039 $\pm$ 44 & 71.9475 $\pm$ 44.1 &  & $\checkmark$ \\
uT5\_lEta2 & 32.5546 $\pm$ 44.9 & 32.464 $\pm$ 44.7 & 32.3097 $\pm$ 44.5 &  & $\checkmark$ \\
uT5\_lEta3 & 30.6633 $\pm$ 49.7 & 30.7196 $\pm$ 49.8 & 30.9355 $\pm$ 50 &  & $\checkmark$ \\
uT5\_lEta4 & 61.7726 $\pm$ 76.7 & 60.9313 $\pm$ 75.7 & 62.0634 $\pm$ 77 &  & $\checkmark$ \\
\hline
lEta1 & 638.32 $\pm$ 116 & 661.456 $\pm$ 118 & 660.775 $\pm$ 118 &  & $\checkmark$ \\
lEta2 & 288.203 $\pm$ 117 & 295.879 $\pm$ 119 & 293.98 $\pm$ 119 &  & $\checkmark$ \\
lEta3 & 444.985 $\pm$ 146 & 457.937 $\pm$ 149 & 462.057 $\pm$ 151 &  & $\checkmark$ \\
lEta4 & 498.806 $\pm$ 223 & 505.058 $\pm$ 225 & 505.705 $\pm$ 225 &  & $\checkmark$ \\
\hline
 &  &  &  & elPt=1883.1 $\pm$ 282 &  \\
 &  &  &  & mT=1941.89 $\pm$ 288 &  \\
 &  &  &  & ptn=1941.58 $\pm$ 288 &  \\
\hline
\end{tabular}
\end{table}
\begin{table}[htbp]
\centering
\scriptsize
\caption{SR MJ yield for wminusenu5TeV}
\label{tab:wminusenu5TeV}
\setlength{\tabcolsep}{6pt}
\begin{tabular}{lrrrrr}
\hline
bin & ptl & mT & ptn & Inclusive & Consistency \\
\hline
uT1\_lEta1 & 54.4909 $\pm$ 57.2 & 13.8796 $\pm$ 22.7 & 43.9284 $\pm$ 27.2 &  & $\checkmark$ \\
uT1\_lEta2 & 82.2254 $\pm$ 65 & 25.7083 $\pm$ 26.7 & 31.7239 $\pm$ 26.6 &  & $\checkmark$ \\
uT1\_lEta3 & 16.4749 $\pm$ 33.3 & 26.4253 $\pm$ 31.9 & 8.95516 $\pm$ 16.5 &  & $\checkmark$ \\
uT1\_lEta4 & 69.0895 $\pm$ 55.9 & 62.4818 $\pm$ 48 & 71.3323 $\pm$ 51.8 &  & $\checkmark$ \\
uT2\_lEta1 & 164.192 $\pm$ 48.5 & 178.116 $\pm$ 46.8 & 167.658 $\pm$ 44.1 &  & $\checkmark$ \\
uT2\_lEta2 & 143.609 $\pm$ 49.6 & 155.437 $\pm$ 48.3 & 152.284 $\pm$ 47.4 &  & $\checkmark$ \\
uT2\_lEta3 & 135.415 $\pm$ 52.1 & 146.715 $\pm$ 54.5 & 147.953 $\pm$ 54.9 &  & $\checkmark$ \\
uT2\_lEta4 & 55.458 $\pm$ 67.1 & 64.1049 $\pm$ 74.7 & 65.816 $\pm$ 76.7 &  & $\checkmark$ \\
uT3\_lEta1 & 106.405 $\pm$ 43.8 & 115.317 $\pm$ 46.8 & 116.024 $\pm$ 47.1 &  & $\checkmark$ \\
uT3\_lEta2 & 178.407 $\pm$ 52.2 & 187.591 $\pm$ 54.4 & 191.818 $\pm$ 55.6 &  & $\checkmark$ \\
uT3\_lEta3 & 119.604 $\pm$ 60.8 & 114.402 $\pm$ 58.1 & 113.639 $\pm$ 57.7 &  & $\checkmark$ \\
uT3\_lEta4 & 191.484 $\pm$ 89.4 & 198.002 $\pm$ 92.3 & 202.017 $\pm$ 94.2 &  & $\checkmark$ \\
uT4\_lEta1 & 78.2779 $\pm$ 41.2 & 78.6245 $\pm$ 41.3 & 79.8047 $\pm$ 41.9 &  & $\checkmark$ \\
uT4\_lEta2 & 64.7903 $\pm$ 53.1 & 64.2675 $\pm$ 52.6 & 64.4273 $\pm$ 52.7 &  & $\checkmark$ \\
uT4\_lEta3 & 69.4727 $\pm$ 45.8 & 71.7129 $\pm$ 47.2 & 71.9702 $\pm$ 47.4 &  & $\checkmark$ \\
uT4\_lEta4 & 183.724 $\pm$ 77.6 & 194.787 $\pm$ 82.2 & 196.995 $\pm$ 83.1 &  & $\checkmark$ \\
uT5\_lEta1 & 57.6504 $\pm$ 35.2 & 57.6404 $\pm$ 35.2 & 58.3444 $\pm$ 35.6 &  & $\checkmark$ \\
uT5\_lEta2 & 69.9919 $\pm$ 42.6 & 69.7604 $\pm$ 42.4 & 69.6815 $\pm$ 42.3 &  & $\checkmark$ \\
uT5\_lEta3 & 28.7812 $\pm$ 40.8 & 28.8224 $\pm$ 40.8 & 28.7826 $\pm$ 40.8 &  & $\checkmark$ \\
uT5\_lEta4 & 23.3316 $\pm$ 68.9 & 23.5576 $\pm$ 69.5 & 24.3628 $\pm$ 72 &  & $\checkmark$ \\
\hline
lEta1 & 510.641 $\pm$ 102 & 528.988 $\pm$ 104 & 532.886 $\pm$ 105 &  & $\checkmark$ \\
lEta2 & 512.48 $\pm$ 122 & 516.94 $\pm$ 122 & 517.23 $\pm$ 122 &  & $\checkmark$ \\
lEta3 & 286.702 $\pm$ 126 & 290.582 $\pm$ 128 & 289.296 $\pm$ 127 &  & $\checkmark$ \\
lEta4 & 362.503 $\pm$ 196 & 378.075 $\pm$ 204 & 385.675 $\pm$ 208 &  & $\checkmark$ \\
\hline
 &  &  &  & elPt=1778.68 $\pm$ 258 &  \\
 &  &  &  & mT=1823.9 $\pm$ 263 &  \\
 &  &  &  & ptn=1830.18 $\pm$ 264 &  \\
\hline
\end{tabular}
\end{table}
\begin{table}[htbp]
\centering
\scriptsize
\caption{SR MJ yield for wplusmunu13TeV}
\label{tab:wplusmunu13TeV}
\setlength{\tabcolsep}{6pt}
\begin{tabular}{lrrrrr}
\hline
bin & ptl & mT & ptn & Inclusive & Consistency \\
\hline
uT1\_lEta1 & 288.752 $\pm$ 138 & 258.124 $\pm$ 94.3 & 268.466 $\pm$ 90.8 &  & $\checkmark$ \\
uT1\_lEta2 & 165.452 $\pm$ 140 & 154.784 $\pm$ 96.4 & 119.147 $\pm$ 90.1 &  & $\checkmark$ \\
uT1\_lEta3 & 205.049 $\pm$ 158 & 152.737 $\pm$ 109 & 23.3899 $\pm$ 101 &  & x \\
uT1\_lEta4 & 1.50489e-05 $\pm$ 51.2 & 12.6395 $\pm$ 97.2 & -0.00097627 $\pm$ 44.6 &  & $\checkmark$ \\
uT2\_lEta1 & 494.999 $\pm$ 104 & 605.868 $\pm$ 84.1 & 597.626 $\pm$ 83.2 &  & x \\
uT2\_lEta2 & 379.468 $\pm$ 103 & 425.535 $\pm$ 75.3 & 417.709 $\pm$ 74.4 &  & $\checkmark$ \\
uT2\_lEta3 & 357.895 $\pm$ 107 & 408.506 $\pm$ 77.7 & 443.634 $\pm$ 80.8 &  & $\checkmark$ \\
uT2\_lEta4 & 88.7586 $\pm$ 71 & 151.611 $\pm$ 52.7 & 138.861 $\pm$ 50.4 &  & $\checkmark$ \\
uT3\_lEta1 & 619.914 $\pm$ 84 & 624.609 $\pm$ 77 & 621.682 $\pm$ 76.8 &  & $\checkmark$ \\
uT3\_lEta2 & 320.531 $\pm$ 63.9 & 330.433 $\pm$ 57.4 & 327.615 $\pm$ 57.1 &  & $\checkmark$ \\
uT3\_lEta3 & 405.596 $\pm$ 74.4 & 430.955 $\pm$ 68.3 & 423.193 $\pm$ 67.5 &  & $\checkmark$ \\
uT3\_lEta4 & 122.083 $\pm$ 52.9 & 121.233 $\pm$ 46.7 & 117.565 $\pm$ 45.6 &  & $\checkmark$ \\
uT4\_lEta1 & 307.298 $\pm$ 70.4 & 299.72 $\pm$ 67.1 & 300.045 $\pm$ 67.2 &  & $\checkmark$ \\
uT4\_lEta2 & 248.866 $\pm$ 58.3 & 246.137 $\pm$ 55.9 & 245.364 $\pm$ 55.8 &  & $\checkmark$ \\
uT4\_lEta3 & 314.563 $\pm$ 65.7 & 309.895 $\pm$ 62.9 & 307.574 $\pm$ 62.5 &  & $\checkmark$ \\
uT4\_lEta4 & 91.6283 $\pm$ 44.6 & 88.1641 $\pm$ 41.6 & 87.6383 $\pm$ 41.3 &  & $\checkmark$ \\
uT5\_lEta1 & 291.484 $\pm$ 63.5 & 285.843 $\pm$ 61.8 & 285.523 $\pm$ 61.7 &  & $\checkmark$ \\
uT5\_lEta2 & 157.083 $\pm$ 51.3 & 150.87 $\pm$ 48.8 & 150.153 $\pm$ 48.5 &  & $\checkmark$ \\
uT5\_lEta3 & 164.085 $\pm$ 55.2 & 162.058 $\pm$ 54 & 161.912 $\pm$ 54 &  & $\checkmark$ \\
uT5\_lEta4 & 104.105 $\pm$ 39.3 & 103.037 $\pm$ 38.4 & 103.472 $\pm$ 38.5 &  & $\checkmark$ \\
\hline
lEta1 & 1976.04 $\pm$ 170 & 1980.79 $\pm$ 158 & 1972.35 $\pm$ 157 &  & $\checkmark$ \\
lEta2 & 1322.73 $\pm$ 140 & 1312.33 $\pm$ 126 & 1297 $\pm$ 125 &  & $\checkmark$ \\
lEta3 & 1517.9 $\pm$ 155 & 1539.68 $\pm$ 142 & 1522.85 $\pm$ 141 &  & $\checkmark$ \\
lEta4 & 506.455 $\pm$ 105 & 527.862 $\pm$ 95.8 & 523.871 $\pm$ 95.4 &  & $\checkmark$ \\
\hline
 &  &  &  & mT=5345.25 $\pm$ 265 &  \\
 &  &  &  & muPt=5314.17 $\pm$ 290 &  \\
 &  &  &  & ptn=5299.65 $\pm$ 263 &  \\
\hline
\end{tabular}
\end{table}
\begin{table}[htbp]
\centering
\scriptsize
\caption{SR MJ yield for wminusmunu13TeV}
\label{tab:wminusmunu13TeV}
\setlength{\tabcolsep}{6pt}
\begin{tabular}{lrrrrr}
\hline
bin & ptl & mT & ptn & Inclusive & Consistency \\
\hline
uT1\_lEta1 & 226.67 $\pm$ 121 & 267.983 $\pm$ 96 & 239.253 $\pm$ 95.3 &  & $\checkmark$ \\
uT1\_lEta2 & 223.794 $\pm$ 102 & 155.067 $\pm$ 80.4 & 168.097 $\pm$ 79.1 &  & $\checkmark$ \\
uT1\_lEta3 & 197.018 $\pm$ 110 & 174.302 $\pm$ 85.5 & 155.848 $\pm$ 85.1 &  & $\checkmark$ \\
uT1\_lEta4 & 134.854 $\pm$ 95 & 80.5779 $\pm$ 69.7 & 150.646 $\pm$ 74.7 &  & $\checkmark$ \\
uT2\_lEta1 & 575.286 $\pm$ 97.4 & 630.124 $\pm$ 78.9 & 636.962 $\pm$ 79.9 &  & $\checkmark$ \\
uT2\_lEta2 & 395.628 $\pm$ 82.4 & 470.37 $\pm$ 68.8 & 474.543 $\pm$ 70.2 &  & $\checkmark$ \\
uT2\_lEta3 & 374.969 $\pm$ 86.8 & 399.537 $\pm$ 71.7 & 400.099 $\pm$ 72.1 &  & $\checkmark$ \\
uT2\_lEta4 & 237.642 $\pm$ 73.7 & 225.907 $\pm$ 58.3 & 204.684 $\pm$ 55.2 &  & $\checkmark$ \\
uT3\_lEta1 & 529.89 $\pm$ 73.8 & 541.725 $\pm$ 69.8 & 541.076 $\pm$ 69.9 &  & $\checkmark$ \\
uT3\_lEta2 & 285.169 $\pm$ 59.5 & 286.008 $\pm$ 54.7 & 285.401 $\pm$ 54.8 &  & $\checkmark$ \\
uT3\_lEta3 & 433.918 $\pm$ 74.9 & 412.482 $\pm$ 66.4 & 400.304 $\pm$ 64.8 &  & $\checkmark$ \\
uT3\_lEta4 & 137.528 $\pm$ 45.8 & 138.018 $\pm$ 42.1 & 133.738 $\pm$ 41.2 &  & $\checkmark$ \\
uT4\_lEta1 & 274.997 $\pm$ 65.3 & 273.273 $\pm$ 63.8 & 272.144 $\pm$ 63.5 &  & $\checkmark$ \\
uT4\_lEta2 & 162.957 $\pm$ 57.3 & 156.184 $\pm$ 53.9 & 155.546 $\pm$ 53.8 &  & $\checkmark$ \\
uT4\_lEta3 & 236.738 $\pm$ 59.8 & 233.594 $\pm$ 57.8 & 232.057 $\pm$ 57.5 &  & $\checkmark$ \\
uT4\_lEta4 & 90.5104 $\pm$ 39.7 & 90.316 $\pm$ 38.9 & 89.3838 $\pm$ 38.6 &  & $\checkmark$ \\
uT5\_lEta1 & 171.469 $\pm$ 56.2 & 169.169 $\pm$ 55.1 & 169.984 $\pm$ 55.4 &  & $\checkmark$ \\
uT5\_lEta2 & 127.618 $\pm$ 49 & 124.304 $\pm$ 47.3 & 123.397 $\pm$ 47 &  & $\checkmark$ \\
uT5\_lEta3 & 151.976 $\pm$ 53.4 & 145.029 $\pm$ 50.7 & 145.265 $\pm$ 50.8 &  & $\checkmark$ \\
uT5\_lEta4 & 42.0157 $\pm$ 35.1 & 41.7483 $\pm$ 34.6 & 41.4913 $\pm$ 34.4 &  & $\checkmark$ \\
\hline
lEta1 & 1867.45 $\pm$ 155 & 1886.27 $\pm$ 146 & 1884.86 $\pm$ 147 &  & $\checkmark$ \\
lEta2 & 1228.05 $\pm$ 132 & 1218.64 $\pm$ 122 & 1207.96 $\pm$ 121 &  & $\checkmark$ \\
lEta3 & 1270.12 $\pm$ 147 & 1214.22 $\pm$ 133 & 1199.94 $\pm$ 132 &  & $\checkmark$ \\
lEta4 & 508.694 $\pm$ 99.5 & 490.114 $\pm$ 90.1 & 480.455 $\pm$ 88.7 &  & $\checkmark$ \\
\hline
 &  &  &  & mT=4805.32 $\pm$ 249 &  \\
 &  &  &  & muPt=4888.87 $\pm$ 269 &  \\
 &  &  &  & ptn=4770.06 $\pm$ 247 &  \\
\hline
\end{tabular}
\end{table}
\begin{table}[htbp]
\centering
\scriptsize
\caption{SR MJ yield for wplusenu13TeV}
\label{tab:wplusenu13TeV}
\setlength{\tabcolsep}{6pt}
\begin{tabular}{lrrrrr}
\hline
bin & ptl & mT & ptn & Inclusive & Consistency \\
\hline
uT1\_lEta1 & 619.536 $\pm$ 152 & 747.664 $\pm$ 139 & 717.965 $\pm$ 135 &  & $\checkmark$ \\
uT1\_lEta2 & 796.128 $\pm$ 172 & 852.091 $\pm$ 150 & 811.396 $\pm$ 145 &  & $\checkmark$ \\
uT1\_lEta3 & 928.007 $\pm$ 171 & 694.473 $\pm$ 128 & 670.195 $\pm$ 123 &  & x \\
uT1\_lEta4 & 860.574 $\pm$ 205 & 711.849 $\pm$ 157 & 604.604 $\pm$ 141 &  & x \\
uT2\_lEta1 & 1186.07 $\pm$ 142 & 1326.41 $\pm$ 136 & 1313.63 $\pm$ 134 &  & $\checkmark$ \\
uT2\_lEta2 & 1250.46 $\pm$ 140 & 1401.65 $\pm$ 135 & 1434.88 $\pm$ 138 &  & x \\
uT2\_lEta3 & 1655.89 $\pm$ 176 & 1674.32 $\pm$ 167 & 1684.8 $\pm$ 168 &  & $\checkmark$ \\
uT2\_lEta4 & 1377.18 $\pm$ 202 & 1405.03 $\pm$ 196 & 1403.1 $\pm$ 195 &  & $\checkmark$ \\
uT3\_lEta1 & 1062.24 $\pm$ 127 & 1090.66 $\pm$ 126 & 1094.32 $\pm$ 127 &  & $\checkmark$ \\
uT3\_lEta2 & 1432.54 $\pm$ 146 & 1463.17 $\pm$ 146 & 1467.19 $\pm$ 146 &  & $\checkmark$ \\
uT3\_lEta3 & 1861.04 $\pm$ 178 & 1853.85 $\pm$ 174 & 1847.44 $\pm$ 173 &  & $\checkmark$ \\
uT3\_lEta4 & 1503.27 $\pm$ 212 & 1519.11 $\pm$ 212 & 1519.17 $\pm$ 212 &  & $\checkmark$ \\
uT4\_lEta1 & 730.843 $\pm$ 121 & 729.476 $\pm$ 120 & 731.431 $\pm$ 121 &  & $\checkmark$ \\
uT4\_lEta2 & 1160.26 $\pm$ 139 & 1162.96 $\pm$ 139 & 1167.67 $\pm$ 139 &  & $\checkmark$ \\
uT4\_lEta3 & 988.856 $\pm$ 152 & 971.469 $\pm$ 149 & 971.94 $\pm$ 149 &  & $\checkmark$ \\
uT4\_lEta4 & 1158.15 $\pm$ 219 & 1146.79 $\pm$ 217 & 1148.64 $\pm$ 217 &  & $\checkmark$ \\
uT5\_lEta1 & 702.05 $\pm$ 108 & 702.321 $\pm$ 107 & 700.032 $\pm$ 107 &  & $\checkmark$ \\
uT5\_lEta2 & 811.539 $\pm$ 115 & 809.715 $\pm$ 115 & 808.947 $\pm$ 114 &  & $\checkmark$ \\
uT5\_lEta3 & 975.525 $\pm$ 134 & 958.877 $\pm$ 131 & 962.149 $\pm$ 132 &  & $\checkmark$ \\
uT5\_lEta4 & 835.49 $\pm$ 191 & 833.579 $\pm$ 190 & 837.395 $\pm$ 191 &  & $\checkmark$ \\
\hline
lEta1 & 4297.22 $\pm$ 284 & 4415.41 $\pm$ 285 & 4409.78 $\pm$ 285 &  & $\checkmark$ \\
lEta2 & 5508.62 $\pm$ 314 & 5654.77 $\pm$ 316 & 5649.71 $\pm$ 315 &  & $\checkmark$ \\
lEta3 & 6069.37 $\pm$ 371 & 5984.84 $\pm$ 360 & 5978.68 $\pm$ 359 &  & $\checkmark$ \\
lEta4 & 5440.92 $\pm$ 498 & 5466.29 $\pm$ 496 & 5470.37 $\pm$ 496 &  & $\checkmark$ \\
\hline
 &  &  &  & elPt=21405.4 $\pm$ 704 &  \\
 &  &  &  & mT=21740.7 $\pm$ 703 &  \\
 &  &  &  & ptn=21740.8 $\pm$ 703 &  \\
\hline
\end{tabular}
\end{table}
\begin{table}[htbp]
\centering
\scriptsize
\caption{SR MJ yield for wminusenu13TeV}
\label{tab:wminusenu13TeV}
\setlength{\tabcolsep}{6pt}
\begin{tabular}{lrrrrr}
\hline
bin & ptl & mT & ptn & Inclusive & Consistency \\
\hline
uT1\_lEta1 & 487.55 $\pm$ 133 & 573.205 $\pm$ 112 & 600.646 $\pm$ 113 &  & $\checkmark$ \\
uT1\_lEta2 & 623.063 $\pm$ 139 & 754.883 $\pm$ 129 & 753.015 $\pm$ 128 &  & $\checkmark$ \\
uT1\_lEta3 & 889.178 $\pm$ 141 & 970.302 $\pm$ 144 & 898.462 $\pm$ 136 &  & $\checkmark$ \\
uT1\_lEta4 & 595.463 $\pm$ 131 & 656.541 $\pm$ 133 & 638.594 $\pm$ 130 &  & $\checkmark$ \\
uT2\_lEta1 & 1035.17 $\pm$ 125 & 1111.64 $\pm$ 123 & 1139.95 $\pm$ 126 &  & $\checkmark$ \\
uT2\_lEta2 & 1513.18 $\pm$ 149 & 1631.54 $\pm$ 146 & 1632.27 $\pm$ 146 &  & $\checkmark$ \\
uT2\_lEta3 & 1524.17 $\pm$ 166 & 1491.91 $\pm$ 157 & 1487.82 $\pm$ 156 &  & $\checkmark$ \\
uT2\_lEta4 & 1679.59 $\pm$ 202 & 1694.52 $\pm$ 199 & 1690.29 $\pm$ 198 &  & $\checkmark$ \\
uT3\_lEta1 & 1188 $\pm$ 136 & 1200.75 $\pm$ 134 & 1199.56 $\pm$ 134 &  & $\checkmark$ \\
uT3\_lEta2 & 1427.51 $\pm$ 137 & 1452.33 $\pm$ 137 & 1464.31 $\pm$ 139 &  & $\checkmark$ \\
uT3\_lEta3 & 1471.43 $\pm$ 162 & 1461.16 $\pm$ 159 & 1462.59 $\pm$ 159 &  & $\checkmark$ \\
uT3\_lEta4 & 1490.86 $\pm$ 211 & 1506.8 $\pm$ 212 & 1520.33 $\pm$ 214 &  & $\checkmark$ \\
uT4\_lEta1 & 1018.17 $\pm$ 127 & 1009.08 $\pm$ 125 & 1006.1 $\pm$ 125 &  & $\checkmark$ \\
uT4\_lEta2 & 939.301 $\pm$ 132 & 930.34 $\pm$ 130 & 926.252 $\pm$ 130 &  & $\checkmark$ \\
uT4\_lEta3 & 1072.5 $\pm$ 154 & 1044.27 $\pm$ 149 & 1041.67 $\pm$ 149 &  & $\checkmark$ \\
uT4\_lEta4 & 1017.54 $\pm$ 204 & 1007.38 $\pm$ 201 & 1010.48 $\pm$ 202 &  & $\checkmark$ \\
uT5\_lEta1 & 783.514 $\pm$ 112 & 777.951 $\pm$ 111 & 779.143 $\pm$ 112 &  & $\checkmark$ \\
uT5\_lEta2 & 816.617 $\pm$ 121 & 808.699 $\pm$ 119 & 810.006 $\pm$ 120 &  & $\checkmark$ \\
uT5\_lEta3 & 703.091 $\pm$ 127 & 691.416 $\pm$ 124 & 690.876 $\pm$ 124 &  & $\checkmark$ \\
uT5\_lEta4 & 484.801 $\pm$ 176 & 487.984 $\pm$ 177 & 491.733 $\pm$ 179 &  & $\checkmark$ \\
\hline
lEta1 & 4503.52 $\pm$ 289 & 4555.12 $\pm$ 286 & 4550.85 $\pm$ 286 &  & $\checkmark$ \\
lEta2 & 5433.33 $\pm$ 309 & 5489.13 $\pm$ 307 & 5484.27 $\pm$ 307 &  & $\checkmark$ \\
lEta3 & 5391.77 $\pm$ 358 & 5279.5 $\pm$ 347 & 5260.95 $\pm$ 346 &  & $\checkmark$ \\
lEta4 & 5042.95 $\pm$ 468 & 5124.87 $\pm$ 473 & 5155 $\pm$ 476 &  & $\checkmark$ \\
\hline
 &  &  &  & elPt=20687.3 $\pm$ 688 &  \\
 &  &  &  & mT=20791.9 $\pm$ 684 &  \\
 &  &  &  & ptn=20778.6 $\pm$ 683 &  \\
\hline
\end{tabular}
\end{table}

\todo{The git repository contains a full plotbook in the subfolder of the figures. I think for documentation purposes this should be attached somehow.}
